\chapter{DEFINITION DU PROBLEME}
\thispagestyle{empty}

\section{Definition du probleme}

Le projet traite le probleme suivant: etant donne un article source de
Wikipedia dans le domaine de l'informatique, retrouver une liste classee
d'articles lies issus du meme domaine. La sortie doit etre utile comme
recommandation et non seulement comme prediction generique de liens. En
pratique, cela signifie que le systeme doit retrouver des articles proches sur
le plan semantique, relies structurellement et pertinents a faire apparaitre en
tete d'une liste courte de recommandations.

Le pipeline de recommandation est donc une approche d'article a article, fondee
sur une requete. Il ne repose ni sur un historique d'utilisateurs, ni sur du
filtrage collaboratif, ni sur des profils personnalises. L'article source joue
le role d'ancre de requete et l'espace candidat correspond a la collection
Custom-WikiCS.

\section{Entrees, sorties et limites}

Les entrees du systeme sont les suivantes:
\begin{itemize}
    \item un graphe derive de Wiki-CS;
    \item le texte rehydrate des articles Wikipedia;
    \item des embeddings semantiques et structurels precalcules;
    \item un titre d'article selectionne comme source de requete.
\end{itemize}

La sortie principale est une liste top-$k$ classee, typiquement un top-20 dans
l'application web locale et un top-$K$ dans plusieurs metriques d'evaluation
hors ligne. Une sortie secondaire optionnelle consiste en une presentation
structuree de type parcours d'apprentissage, qui reordonne les titres deja
recuperes dans des categories telles que fondements, concepts centraux, sujets
avances et sujets adjacents.

Les limites du projet sont egalement claires:
\begin{itemize}
    \item il est restreint aux articles d'informatique de Wikipedia en anglais;
    \item il ne constitue pas un service web de production;
    \item il ne s'agit pas d'un systeme de recommandation personnalise;
    \item la couche LLM n'intervient qu'en post-traitement et n'est pas
    consideree comme le moteur central de recuperation.
\end{itemize}

\section{Contraintes realistes}

\subsection{Calcul et economie}

Toutes les experiences ont ete menees dans un cadre academique local plutot que
sur une infrastructure cloud. Cette contrainte a fortement influence la
conception:
\begin{itemize}
    \item l'entrainement et l'evaluation ont ete realises sur un ordinateur
    portable grand public equipe d'un RTX 4070 Laptop GPU et de 8GB de VRAM;
    \item les artefacts d'embeddings ont ete mis en cache afin d'eviter des
    calculs couteux repetes;
    \item l'organisation du travail privilegie les notebooks et
    l'experimentation locale plutot qu'un entrainement distribue.
\end{itemize}

Cette conception orientee local a rendu le projet economique et reproductible,
mais elle a egalement limite l'echelle et la complexite des experiences.

\subsection{Passage a l'echelle}

Le graphe experimental nettoye contient 11,367 noeuds actifs et 291,039 aretes
orientees. Cette taille est suffisante pour faire apparaitre une structure de
graphe significative tout en restant gerable sur du materiel local. Un
entrainement a l'echelle complete de Wikipedia exigerait d'autres choix
d'ingenierie, tels que l'echantillonnage sur graphe, une infrastructure plus
importante ou un traitement distribue. Le projet privilegie donc la correction,
la clarte et la reproductibilite locale plutot qu'un deploiement a l'echelle
du web.

\subsection{Ethique, confidentialite et securite}

Le projet ne traite aucune donnee personnelle d'utilisateur. La recommandation
est fondee sur la requete, sans suivi comportemental ni profilage. Cela rend le
systeme plus facile a auditer et reduit le risque pour la vie privee. La
logique centrale de classement reste egalement transparente: elle repose sur
des embeddings fixes en cache et sur des calculs explicites de similarite.

La couche LLM introduit un autre type de risque: l'inconsistance de sortie.
Pour la reduire, l'organisateur est contraint d'operer uniquement sur des
titres deja recuperes et sa sortie est validee avant presentation. Si la
reponse structuree est inutilisable, le systeme revient a un parcours
heuristique deterministe.

\subsection{Validite scientifique}

L'une des contraintes les plus importantes de ce projet concerne la validite
methodologique. Comme les modeles structurels peuvent exploiter l'information du
graphe de facon subtile, la definition des decoupages apprentissage/test est
determinante. Les premieres experiences combinaient une retenue orientee avec
une modelisation structurelle qui consommait en pratique un graphe non oriente.
Cette incoherence a motive un protocole d'evaluation revise, devenu central
dans le rapport final.

\section{Criteres de succes}

Le projet s'appuie sur plusieurs niveaux de criteres de succes:
\begin{itemize}
    \item la structure du graphe doit rester significative apres enrichissement
    et nettoyage;
    \item les embeddings semantiques et structurels doivent tous deux apporter
    un signal exploitable pour la recommandation;
    \item le protocole d'evaluation doit eviter toute incoherence
    methodologique;
    \item la qualite finale du classement doit etre mesuree a la fois par des
    metriques de recuperation et par des metriques de qualite de recommandation;
    \item le prototype web doit exposer la famille de modeles dans une
    interface locale utilisable.
\end{itemize}

Ce cadrage multicritere est important, car le meilleur modele selon une famille
de metriques n'est pas automatiquement le meilleur pour l'objectif reel du
produit, a savoir une recommandation d'articles en liste courte et classee.
